% 四类填充示例共用的正文；应由同目录的 xdupgtm-*.tex 引入。
% 本文件中的姓名、项目、实验结果和成果均用于排版测试，不可直接提交。
% signature-sample.png 和 date-sample.png 来自 note286/xdupgtp 的公开示例。
% 所有签名位置重复使用同一示例图片，仅检验尺寸、换行和单元格边界。

% 基本信息示例：长题目分成两行，测试封面下划线和日期。
\newcommand\studentNo{1101110071}
\newcommand\thesisTitleOne{复杂电磁环境下的高效建模与实验验证}
\newcommand\thesisTitleTwo{中期考核报告（填充测试示例）}
\newcommand\authorName{张三}
\ifadegree
  \newcommand\major{电子科学与技术}
  \newcommand\submajor{电磁场与微波技术}
\else
  \newcommand\major{电子信息}
  \newcommand\submajor{新一代电子信息技术}
\fi
\newcommand\advisorName{李四}
\newcommand\entadvisorName{王五}
\newcommand\school{电子工程学院}
\newcommand\submitDate{2026年10月9日}
\newcommand\subjectSource{科研项目与工程应用需求（虚构测试项目）}

% 学术学位导师、专业学位校内导师及组长：图片签名配合文字日期。
\renewcommand\zdjsqm{\sign{examples/filled/signature-sample.png}}
\renewcommand\zdjsqmrq{2026年10月9日}
\renewcommand\xnjsqm{\sign{examples/filled/signature-sample.png}}
\renewcommand\xnjsqmrq{2026年10月9日}
\renewcommand\zzqm{\sign{examples/filled/signature-sample.png}}
% 校外导师：图片签名配合图片日期，测试不同宽高比的图片。
\renewcommand\xyjsqm{\sign{examples/filled/signature-sample.png}}
\renewcommand\xyjsqmrq{\sign{examples/filled/date-sample.png}}
% 成员签名：emptysign 在相邻签名之间预留空白，超宽时应自动换行。
\renewcommand\cyqm{\sign{examples/filled/signature-sample.png}\emptysign
  \sign{examples/filled/signature-sample.png}\emptysign
  \sign{examples/filled/signature-sample.png}}
\renewcommand\cyqmrq{2026年10月9日}

% 自动化测试可定义 xdupgtmStressTest，增加长正文和五名成员签名。
% 普通填充示例不启用此模式；执行 test-filled.ps1 可同时生成两种版本。
\newcount\exampleParagraphCount
\newcommand\exampleExtraParagraphs[1]{%
  \ifdefined\xdupgtmStressTest
    \exampleParagraphCount=0
    \loop\ifnum\exampleParagraphCount<#1
      \advance\exampleParagraphCount by 1
      \par 扩展验证段落（第\the\exampleParagraphCount 段）：本阶段继续考察不同边界条件、网格密度与测量噪声对模型输出的影响。全部算例采用相同的数据划分和评价指标，记录计算时间、相对误差及重复实验的波动范围。对于偏差较大的样本，逐项核查材料参数、激励设置和校准过程，并保存输入文件、程序版本以及分析记录。补充实验按照统一流程执行，使用 MATLAB 和 Python 交叉核对结果，为后续论文中的误差分析和工程验证提供依据。\par
    \repeat
  \fi}
\ifdefined\xdupgtmStressTest
  \renewcommand\cyqm{\sign{examples/filled/signature-sample.png}\emptysign
    \sign{examples/filled/signature-sample.png}\emptysign
    \sign{examples/filled/signature-sample.png}\emptysign
    \sign{examples/filled/signature-sample.png}\emptysign
    \sign{examples/filled/signature-sample.png}}
\fi

\begin{document}
% 选题来源通过上面的 subjectSource 填写，来源行与第一部分边框相连。
\begin{xzlyj}\end{xzlyj}
\begin{yjmbnr}
% 第一部分示例：使用正常段落、英文术语及编号列表描述目标和内容。
本课题面向复杂电磁环境中的快速预测与可靠验证需求，研究兼顾计算效率与精度的建模方法。研究目标是在保持物理约束的基础上，建立可复用的数值计算流程，并通过仿真算例和实验测量检验其适用范围。研究目标及研究内容与开题报告保持一致。

拟开展的研究内容包括：分析结构参数与电磁响应之间的关系，建立基准模型及误差评价方法；研究参数化模型的降阶计算策略，降低重复计算开销；搭建验证平台，对典型工况进行测量并分析模型偏差。

\begin{enumerate}
  \item 建立可追溯的参数与数据管理流程，统一输入格式、单位和评价指标。
  \item 对比 Full-wave simulation 与降阶模型，评估精度、速度和稳定性。
  \item 采用独立测试样本检验模型的泛化能力，并明确方法的应用边界。
\end{enumerate}
研究中将保留原始数据、计算脚本和处理日志。所有测试样本均分离于模型拟合过程，避免训练数据与验证数据混用。
\exampleExtraParagraphs{14}
本部分的目标和研究内容已完整列出。
\end{yjmbnr}

\begin{ywcgz}
% 第二部分示例：段落与小型表格混排，数字和西文使用 Times New Roman。
目前已完成文献整理、基准模型构建及第一轮仿真验证。理论研究方面，梳理了频域求解、模型降阶及不确定性分析方法，确定以相对误差和计算耗时作为主要评价指标，并建立参数归一化与结果复核流程。

数值研究方面，完成三类典型结构的网格划分与边界条件设置，构建了包含 120 组参数配置的测试数据。对相同配置重复计算后，输出结果保持稳定。对部分差异较大的算例进行了网格加密验证，确认误差主要来自结构细节和材料参数变化。

\begin{center}
\begin{tabular}{lrr}
\toprule
验证工况 & 相对误差（\%） & 耗时（s）\\
\midrule
基准工况 A & 1.8 & 12.5\\
扰动工况 B & 2.4 & 14.1\\
独立工况 C & 3.1 & 15.7\\
\bottomrule
\end{tabular}
\end{center}
% 表内数据均为虚构，不能作为学术结果引用。
上述虚构数据仅用于测试表格排版。实验方面，已完成设备校准和初步测量，建立了包含样本编号、环境条件及重复次数的记录表。已获得的阶段性结论是：在当前验证范围内，参数化模型能够保持较小偏差，但对复杂边界条件仍需补充实验。

本阶段形成了基准算例、数据处理脚本和实验记录三类材料。所有材料按版本归档，下一阶段将增加独立样本并开展误差来源分析。
\exampleExtraParagraphs{16}
本部分的已完成工作及阶段结论已完整列出。
\end{ywcgz}

\begin{workplan}
% 第三至第五部分合用一个表格，长内容应在此表格内续页。
\begin{bxfsm}
% 没有偏离时可以直接填写“无”；此处测试包含一项合理调整的说明。
研究目标与开题报告一致。由于实验设备排期调整，原计划在本阶段完成的高频测量顺延至下一阶段，已通过补充仿真验证先行核查主要假设。该调整不改变研究范围和预期成果，后续将按计划补齐实验数据。
\exampleExtraParagraphs{3}
本部分的研究计划调整情况已完整说明。
\end{bxfsm}
\begin{gzhjh}
% 下一步计划可按时间列出，使用列表测试缩进及表格内换行。
下一阶段拟依次完成以下工作：
\begin{enumerate}
  \item 2026年10月至11月：完善复杂边界工况，补充网格收敛与参数敏感性分析。
  \item 2026年12月至2027年1月：开展独立测量，核查实验与仿真偏差，完成交叉验证。
  \item 2027年2月至3月：整理阶段成果，完成论文初稿并根据导师意见修改。
\end{enumerate}
需要解决的关键技术包括复杂结构的稳定参数化、有限样本下的误差评估，以及测量数据与仿真结果之间的统一校准。
\exampleExtraParagraphs{16}
本部分的后续计划和关键技术已完整列出。
\end{gzhjh}
\begin{kycg}
% 区分已发表、已录用和拟发表，避免将计划成果写成已完成成果。
已形成基准模型、计算脚本及阶段实验记录，相关材料均为本示例的虚构内容。截至本次考核，尚无已发表或已录用论文。

拟在完成独立验证后整理一篇关于参数化建模与误差分析的论文，目前处于材料准备阶段。另计划将数据整理流程形成可复现的研究附件，随论文初稿提交导师审阅。
\exampleExtraParagraphs{10}
本部分的成果及发表状态已完整列出。
\end{kycg}
\end{workplan}

% 导师意见保持简洁，以适应官方表格的固定高度。
\ifadegree
\begin{jsyj}
该生已完成基准模型和初步验证，研究进度符合计划。建议补充独立实验并加强误差分析，继续推进论文工作。
\end{jsyj}
\else
\begin{xnjsyj}
该生已完成基准模型和初步验证，进度符合计划。建议补充独立实验，加强误差分析。
\end{xnjsyj}
\begin{xyjsyj}
研究内容与工程需求一致，阶段材料完整。建议补充实际工况测量并完善验证记录。
\end{xyjsyj}
\fi

\begin{pyjl}
% 此处为虚构评语；正式报告中的结论和签名由考核人员确认。
考核组认为该生已完成预定的阶段研究任务，报告结构完整，能够说明研究进展和后续计划。建议增加独立样本，进一步分析误差来源，并明确方法的适用边界。中期考核结论：通过，按专家意见修改后继续学位论文撰写工作。
\end{pyjl}
\end{document}
